\section{归一化层}

\subsection{归一化的基本原理}

归一化层的核心思想分两步：
\begin{enumerate}
    \item \textbf{标准化}：将输入数据变换为均值为 0、标准差为 1 的分布
    \item \textbf{可学习的缩放与平移}：通过可训练参数 $\gamma$（缩放）和 $\beta$（平移），让网络自己决定每一层的最优分布
\end{enumerate}

所有归一化层都遵循同一个公式：

$$\hat{x} = \frac{x - \mu}{\sigma + \epsilon}, \quad y = \gamma \hat{x} + \beta$$

\begin{itemize}
    \item $\mu$, $\sigma$：计算得到的均值和标准差
    \item $\gamma$, $\beta$：可学习参数
    \item $\epsilon$：防止除零的小常数
\end{itemize}

不同归一化方法的\textbf{唯一区别}在于：\textbf{在哪些维度上计算 $\mu$ 和 $\sigma$}。

\subsection{Layer Normalization}

\begin{figure}[H]
\centering
\includegraphics[width=0.95\textwidth]{slides-images/slide-005.jpg}
\caption{Layer Norm：对每个样本的所有通道和空间维度计算统计量\protect\footnotemark}
\end{figure}
\footnotetext{Slides 第6页。}

\textbf{Layer Norm} 是目前最常用的归一化方法（尤其在 Transformer 中）：
\begin{itemize}
    \item 对\textbf{每个样本}独立计算统计量
    \item 计算维度：通道 $\times$ 高度 $\times$ 宽度（对于 CNN）或特征维度 $D$（对于 MLP）
    \item 每个样本只有一个 $\mu$ 和一个 $\sigma$
\end{itemize}

对于向量输入 $x \in \mathbb{R}^{N \times D}$：
$$\mu_n = \frac{1}{D} \sum_{d=1}^{D} x_{n,d}, \quad \sigma_n = \sqrt{\frac{1}{D} \sum_{d=1}^{D} (x_{n,d} - \mu_n)^2}$$

\subsection{Batch Normalization}

\begin{figure}[H]
\centering
\includegraphics[width=0.95\textwidth]{slides-images/slide-007.jpg}
\caption{四种归一化方法对比：Batch Norm、Layer Norm、Instance Norm、Group Norm\protect\footnotemark}
\end{figure}
\footnotetext{Slides 第8页。}

\textbf{Batch Norm} 在每个通道内、跨 batch 中所有样本计算统计量：
\begin{itemize}
    \item 对\textbf{每个通道}独立计算 $\mu$ 和 $\sigma$
    \item 计算维度：batch $\times$ 高度 $\times$ 宽度
    \item 统计量依赖 mini-batch，因此\textbf{训练和推理行为不同}
\end{itemize}

\begin{warningbox}{Batch Norm 的训练/推理不一致}
训练时，Batch Norm 使用当前 mini-batch 的统计量。推理时，使用训练过程中累积的全局均值和方差（running statistics）。这种不一致会带来问题：
\begin{itemize}
    \item 小 batch size 时统计量不稳定
    \item 分布式训练时需要跨 GPU 同步统计量
    \item 推理行为依赖训练时的 batch 组成
\end{itemize}
这也是 Layer Norm 在现代架构中逐渐取代 Batch Norm 的原因之一。
\end{warningbox}

\subsection{Instance Norm 与 Group Norm}

\begin{itemize}
    \item \textbf{Instance Norm}：对每个样本的每个通道独立计算统计量（只在空间维度上归一化）。常用于风格迁移。
    \item \textbf{Group Norm}：将通道分组，在每组内计算统计量。是 Layer Norm 和 Instance Norm 的折中。
\end{itemize}

\begin{importantbox}{归一化方法选择指南}
\begin{itemize}
    \item \textbf{Transformer / 序列模型}：Layer Norm（标准选择）
    \item \textbf{传统 CNN}：Batch Norm（历史上最常用）
    \item \textbf{小 batch / 分布式训练}：Group Norm 或 Layer Norm
    \item \textbf{风格迁移}：Instance Norm
\end{itemize}
\end{importantbox}

\subsection{本章小结}

归一化层通过标准化中间激活值来稳定训练过程。所有归一化方法都遵循``标准化 + 可学习缩放平移''的两步范式，区别仅在于统计量的计算维度。Layer Norm 因其与 batch 无关的特性，成为现代深度学习中最受欢迎的选择。

%% ========== Part 3: Dropout ==========

